\scanpage{107}
\uline{Proposition} 2 Les $a_\alpha$ ($1\leq\alpha\leq3$) sont caractérisés par la condition que $p^3+a_1p^2+a_2p+a_3\mid J=0$.

C'est nécessaire, puisque $p'\mid J=0$. Mais d'après le cor. à la prop. \unsure{précédent|précédente}, les fonctions $1,p,p^2$ restreintes à $J$ sont lin$^{\mathrm t}$ indépendantes, donc la relation précédente sur les $a_i$ (qui s'écrit
\[
 (-p^3)\mid J=a_1(p^2\mid J)+a_2(p\mid J)+a_3(1_E\mid J)\qquad)
\]
caractérise déjà les $a_i$, cqfd.

\uline{Corollaire 1} Les $\unsure{(-1)^\alpha/a_\alpha|\inserted{(-1)^\alpha}a_\alpha}$ ($1\leq\alpha\leq3$) sont les fonctions sym. élémentaires en les $p(i)$ ($i\in J$)

i.e. on a
\[
 p^3+a_1p^2+a_2p+a_3=\prod_{i\in J}\bigl(p-p\deleted{\unsure{f|\xi|'|p}}(i)\bigr)\,. \tag*{(29)}
\]
\deleted{\illegible{1}}
\uline{Corollaire 2} Pour que l'on ait $a_1=\unsure{0|g}$, il f. et s. que $\operatorname{Tr}_{J/k}(p\mid J)=0$ i.e. $\sum p(i)=0$.

Si car. $k\ne3$, les $p\in F_2$ qui satisfont cette condition forment un supplémentaire $W_E$ de $F_1$ dans $F_2$
\[
 F_2\simeq\underbrace{F_1}_{k.1_E}\oplus\underbrace{W_E}_{\simeq\, t_E^{\otimes2}} \tag*{(30)}
\]
Les \deleted{\unsure{sections|fonctions}} \inserted{éléments} de $W_E$ seront \inserted{donc} les fonctions de Weierstrass,
% archive sheet number: 106
